\clearpage
\fancyhead[L]{\small CARLINI. EFFEMERIDI ASTRONOMICHE.}
\hypertarget{native-579-lower}{}
\begin{center}\small Göttingische gelehrte Anzeigen. Stück 10. Seite 89..93. 1815 Januar 19.\\[3pt]\rule{.72\linewidth}{.3pt}\end{center}

{\setlength{\rightskip}{0pt plus 3em}\emph{Effemeridi astronomiche di Milano per l'anno} 1815 \emph{calcolate da} \textsc{Francesco Carlini}.
\emph{Con appendice. Dalla r.\,c. stamperia di governo.} Mailand 1814. Der Kalender 128, der
Anhang 118 Seiten in Octav.\par}

Die stehenden Artikel, welche dem in seiner musterhaften Einrichtung unverändert gebliebenen
astronomischen Kalender beigegeben sind, und die wir bei Gelegenheit der frühern Jahrgänge
angezeigt haben, sind diesmal durch einen sehr schätzbaren Zusatz vermehrt. Dies ist das
Verzeichniss der 34 \textsc{Maskelyne}'schen Sterne aus \textsc{Piazzi}'s neuem Cataloge
entlehnt, zugleich mit der für 1800 und 1850 berechneten Präcession, der eignen Bewegung, und
den Constanten, welche zur Berechnung der Aberration und Nutation nach Hrn. \textsc{von Zach}'s
Manier nöthig sind, so dass man alles vollständig beisammen hat, was zur Bestimmung des
scheinbaren Orts für das ganze gegenwärtige und selbst für das vorhergehende Jahrhundert
erforderlich ist. Der Anhang enthält folgende wichtige Abhandlungen. Breite der Sternwarte Brera
aus den Beobachtungen der Circumpolarsterne von \textsc{Barnabas Oriani}. Nach einer
wohlgeschriebenen Uebersicht der Beobachtungsmethoden, die in der neuern practischen Astronomie
nach und nach bis jetzt zur schärfsten Bestimmung der Polhöhen überhaupt in Anwendung gebracht
worden sind, kommt der vortreffliche Astronom auf die Mailänder Sternwarte insbesondere, und
erzählt die verschiedenen Versuche, ihre Polhöhe immer schärfer zu berichtigen, unter denen die
neuesten von \textsc{Carlini} mit einem 16zolligen \textsc{Lenoir}'schen Repetitionskreise —
demselben, welchen \textsc{Mechain} in Barcellona und Montjoui gebraucht hatte — angestellten
Beobachtungen noch eine Ungewissheit von einigen Secunden zurückliessen. Seit dem Ende des
Jahres 1810 besitzt nun jene Sternwarte einen 3fussigen Kreis mit stehender Säule von
\textsc{Reichenbach}, mit welchem während eines Zeitraums von einem Jahre eine ununterbrochene
Reihe von Beobachtungen gemacht wurden, die in den Jahrgängen der Mailänder Ephemeriden von 1812
und 1813 vollständig abgedruckt sind. Wir erhalten nun gegenwärtig die erste Ausbeute dieses
reichen Schatzes, nemlich die vollständigen Resultate der Beobachtungen vom Polarstern, $\delta$
Cassiopea und $\varepsilon$ im grossen Bär. Die Uebersicht dieser einzelnen Resultate zeigt immer
noch kleine Unterschiede von einem Tage zum andern, die auf mehrere Secunden gehen, und theils
den unvermeidlichen Beobachtungsfehlern, theils Veränderungen in der Atmosphäre, welche durch
Barometer und Thermometer nicht angezeigt werden, zuzuschreiben sein mögen, wenn man nicht einen
Theil davon der Einrichtung des Instruments selbst beimessen will, an welcher die Verbindungsart
des Kreises mit der Säule bekanntlich einige Bedenklichkeiten veranlasst hat. \textsc{Oriani}
wird, der letztern wegen, noch eine bewegliche Libelle am Kreise selbst anbringen lassen, und es
wird höchst interessant sein, die Wirkung dieser Abänderung zu sehen. Der Polarstern wurde an 159
Tagen unter dem Pole, und an 143 Tagen über dem Pole beobachtet. Es ist merkwürdig, dass die
Beobachtungen bei Tage eine schlechtere Uebereinstimmung, zugleich aber eine kleinere
Zenithdistanz geben, als die Beobachtungen bei Nacht. \textsc{Oriani} hat daher jene ganz
ausgeschlossen, und so aus den übrigen abgeleitet
\SourceLeaf{580}
\begin{center}\small\begin{tabular}{@{}l@{\hspace{9pt}}r@{}}
Polhöhe der Sternwarte Brera $.\;.\;.\;.$ & $49^\circ\,28'\,0''713$ \\
Declination des Polarsterns für 1811 $.\;.$ & $88\phantom{\circ}\,17\phantom{'}\,59.494$ \\
\end{tabular}\end{center}
\noindent Nähme man aus \emph{sämmtlichen} Beobachtungen das Mittel, so fände sich nach unsrer Rechnung
\begin{center}\small\begin{tabular}{@{}l@{\hspace{9pt}}r@{}}
Polhöhe der Sternwarte $.\;.\;.\;.\;.\;.\;.$ & $49^\circ\,28'\,1''242$ \\
Declination 1811 $.\;.\;.\;.\;.\;.\;.\;.\;.$ & $88\phantom{\circ}\,17\phantom{'}\,59.583$ \\
\end{tabular}\end{center}
\noindent \textsc{Oriani} hat jenen Weg gewählt, weil er die Unterschiede der Tagbeobachtungen besondern Modificationen der Atmosphäre zuschreibt. Allein ein Theil derselben mag wohl immer auf Rechnung der jährlichen Parallaxe gesetzt werden. Dürften wir von fremden Einflüssen, die am Tage immer oder wenigstens überwiegend in einerlei Sinn wirken, abstrahiren, so würden wir geneigt sein, auf das aus unsrer Rechnung hervorgehende Resultat einer Declinationsparallaxe von $0''7$ im Maximum einiges Gewicht zu legen: es würde dieser eine Rectascensionsparallaxe von $1^{\mathrm{s}}5$ in Zeit entsprechen, welche man um so leichter für zulässig halten könnte, da \textsc{Piazzi} aus den Palermer Beobachtungen sogar das Doppelte gefunden hat. --- Der Stern $\delta$ Cassiopea war in 95 obern und 54 untern, so wie $\varepsilon$ des grossen Bär in 26 obern und 55 untern Culminationen beobachtet, und zwar immer nur bei Nacht oder in der Dämmerung, weshalb \textsc{Oriani} hier keine Beobachtungen auszuschliessen für nöthig fand. Die Polhöhe ergab sich hiernach
\begin{center}\small\begin{tabular}{@{}l@{\hspace{9pt}}r@{}}
aus $\delta$ Cassiopea $.\;.\;.\;.\;.\;.\;.\;.\;.\;.$ & $45^\circ\,28'\,0''975$ \\
aus $\varepsilon$ des grossen Bär $.\;.\;.\;.\;.\;.\;.$ & $45\phantom{\circ}\,28\phantom{'}\,0.\,28$ \\
\end{tabular}\end{center}
\noindent und die Declinationen für 1811
\begin{center}\small\begin{tabular}{@{}l@{\hspace{9pt}}r@{}}
von $\delta$ Cassiopea $.\;.\;.\;.\;.\;.\;.\;.\;.\;.$ & $59^\circ\,14'\,53''415$ \\
von $\varepsilon$ im grossen Bär $.\;.\;.\;.\;.\;.\;.$ & $56\phantom{\circ}\,59\phantom{'}\,15.\,19$ \\
\end{tabular}\end{center}
\noindent Genau genommen liegt indessen in jener nahen Uebereinstimmung nicht sowohl eine Bestätigung der Polhöhe, als ein Beweis für die Güte der bei den Rechnungen angewandten \textsc{Carlini}'schen Refractionstafel, von der wir wünschten, dass sie gleichfalls einen stehenden Artikel der Ephemeriden ausmachen möchte. \textsc{Oriani} fügt am Schluss noch ein Verzeichniss der Declinationen von 3c Circumpolarsternen bei (die obigen drei mitgerechnet), ohne indessen von diesen die Resultate der einzelnen Beobachtungen zu geben. In Rücksicht der Polhöhe stimmen auch die andern Sterne, wie uns \textsc{Oriani} versichert, alle bis auf einige Zehntel einer Secunde mit dem oben gegebenen Resultate durch den Polarstern überein.

Der \emph{zweite} Aufsatz, gleichfalls von \textsc{Oriani}, bestimmt die Opposition des Mars 1813, wo \textsc{Carlini} die geraden Aufsteigungen am Mittagsfernrohr, \textsc{Oriani} die Abweichungen am Kreise beobachtet hatte. Der mittlere Fehler der \textsc{von Lindenau}'schen Marstafeln wurde $-4''8$ in der Länge, $+18''8$ in der Breite, geocentrisch, gefunden. --- Hierauf folgt eine schätzbare Reihe Meridianbeobachtungen der Sonne am Mauerquadranten von \textsc{Angelo Cesaris} von Anfang 1808 bis Ende 1811 als Fortsetzung der in den Jahrgängen 1809 und 1810 gelieferten. --- Der hierauf folgende Aufsatz von \textsc{Carlini} liefert Tafeln für die Mittelpunktsgleichung der Ceres und für die Reduction dieses Planeten auf die Ekliptik. Jene ist für die Excentricität $0.0784$, diese für die Neigung $10^\circ\,37'\,40''$ berechnet; zugleich sind die Aenderungen für $0.0001$ Aenderungen der Excentricität und $10''$ Aenderung der Neigung beigefügt. Aehnliche Tafeln für die Pallas, Juno und Vesta werden für den nächsten Jahrgang versprochen. Wir gestehen, dass wir bei der Pallas und Juno dergleichen Tafeln nicht ganz zweckmässig finden können, insofern die Theorie dieser Planeten auch in Zukunft doch die Form \emph{veränderlicher} Elemente wird behalten müssen, und die Veränderungen viel zu gross sind, um durch solche Tafeln Bequemlichkeit für die Rechnung gewinnen zu können. Um von der Grösse dieser Veränderlichkeit einen Begriff zu geben, bemerken \SourceLeaf{581}wir, dass bei der Opposition des Jahrs 1803 die Excentricität der Pallasbahn $=0.24554$, bei der
Opposition des Jahrs 1814 hingegen $=0.24135$ gewesen ist. Bei der gegenwärtigen Gestalt der
Berechnung der Planetenbewegungen kann man auch dieser Tafeln sehr füglich entbehren. Die
meteorologischen Beobachtungen auf der Mailänder Sternwarte im Jahre 1813 beschliessen den Band:
welch ein glückliches Clima für die praktische Astronomie, wo in Einem Jahre 188 heitere Tage
gezählt werden!
